\documentclass[12pt,a4paper]{article}
\usepackage{../../../myconfig}

\begin{document}

% =========================================================================
% CHUYÊN ĐỀ: PHƯƠNG PHÁP TỌA ĐỘ HÓA
% PHẦN 1: BÀI TẬP RÈN LUYỆN KỸ NĂNG GẮN TRỤC VÀ XÁC ĐỊNH TỌA ĐỘ ĐỈNH
% =========================================================================

% -------------------------------------------------------------------------
% CÂU 1
% -------------------------------------------------------------------------
\titlebox{Chuyên đề: Oxyz}{Phương pháp tọa độ hóa}{Oxyz: Phương pháp tọa độ hóa}

\vspace{-0.2cm}

\questioncore{1}{Cho hình chóp $S.ABC$ có đáy $ABC$ là tam giác vuông tại $A$ với $AB = 2$, $AC = 3$. Cạnh bên $SA$ vuông góc với mặt phẳng đáy và $SA = 4$. Hãy chọn hệ trục tọa độ $Oxyz$ thích hợp, vẽ hình và xác định tọa độ các đỉnh của hình chóp.}

\vspace{0.3cm}
\drawdashlinesfull

\newpage

% -------------------------------------------------------------------------
% CÂU 2
% -------------------------------------------------------------------------
\questioncore{2}{Cho hình chóp $S.ABC$ có đáy $ABC$ là tam giác đều cạnh bằng $2$. Cạnh bên $SA$ vuông góc với mặt phẳng đáy và $SA = 3$. Hãy chọn hệ trục tọa độ $Oxyz$ thích hợp, vẽ hình và xác định tọa độ các đỉnh của hình chóp.}

\vspace{0.3cm}
\drawdashlinesfull

\newpage

% -------------------------------------------------------------------------
% CÂU 3
% -------------------------------------------------------------------------
\questioncore{3}{Cho hình chóp $S.ABC$ có $SA \perp (ABC)$, đáy $ABC$ là tam giác vuông tại $B$ với $AB = 3$, $BC = 4$. Biết chiều cao $SA = 3$. Hãy chọn hệ trục tọa độ $Oxyz$ thích hợp, vẽ hình và xác định tọa độ các đỉnh của hình chóp.}

\vspace{0.3cm}
\drawdashlinesfull

\newpage

% -------------------------------------------------------------------------
% CÂU 4
% -------------------------------------------------------------------------
\questioncore{4}{Cho hình chóp $S.ABCD$ có đáy $ABCD$ là hình chữ nhật với $AB = 3$, $AD = 4$. Cạnh bên $SA$ vuông góc với mặt phẳng đáy và $SA = 5$. Hãy chọn hệ trục tọa độ $Oxyz$ thích hợp, vẽ hình và xác định tọa độ các đỉnh của hình chóp.}

\vspace{0.3cm}
\drawdashlinesfull

\newpage

% -------------------------------------------------------------------------
% CÂU 5
% -------------------------------------------------------------------------
\questioncore{5}{Cho hình chóp $S.ABCD$ có đáy $ABCD$ là hình vuông cạnh bằng $2$. Mặt bên $SAB$ là tam giác đều nằm trong mặt phẳng vuông góc với đáy. Gọi $H$ là trung điểm của cạnh $AB$. Hãy chọn hệ trục tọa độ $Oxyz$ thích hợp, vẽ hình và xác định tọa độ các đỉnh của hình chóp.}

\vspace{0.3cm}
\drawdashlinesfull

\newpage

% -------------------------------------------------------------------------
% CÂU 6
% -------------------------------------------------------------------------
\questioncore{6}{Cho hình tứ diện đều $ABCD$ có tất cả các cạnh bằng $2$. Gọi $H$ là trọng tâm của tam giác $BCD$. Hãy chọn hệ trục tọa độ $Oxyz$ thích hợp, vẽ hình và xác định tọa độ các đỉnh của tứ diện đều.}

\vspace{0.3cm}
\drawdashlinesfull

\newpage

% -------------------------------------------------------------------------
% CÂU 7
% -------------------------------------------------------------------------
\questioncore{7}{Cho hình chóp $S.ABCD$ có đáy $ABCD$ là hình thoi tâm $O$ với $AC = 4$, $BD = 2$. Cạnh bên $SO$ vuông góc với mặt phẳng đáy và $SO = 3$. Hãy chọn hệ trục tọa độ $Oxyz$ thích hợp, vẽ hình và xác định tọa độ các đỉnh của hình chóp.}

\vspace{0.3cm}
\drawdashlinesfull

\newpage

% -------------------------------------------------------------------------
% CÂU 8
% -------------------------------------------------------------------------
\questioncore{8}{Cho hình hộp chữ nhật $ABCD.A'B'C'D'$ có các kích thước $AB = 3$, $AD = 4$, $AA' = 2$. Hãy chọn hệ trục tọa độ $Oxyz$ thích hợp, vẽ hình và xác định tọa độ toàn bộ $8$ đỉnh của hình hộp chữ nhật.}

\vspace{0.3cm}
\drawdashlinesfull

\newpage

% -------------------------------------------------------------------------
% CÂU 9
% -------------------------------------------------------------------------
\questioncore{9}{Cho hình lăng trụ $ABC.A'B'C'$ có đáy $ABC$ là tam giác vuông tại $A$ với $AB = 2$, $AC = 3$. Hình chiếu vuông góc của đỉnh $A'$ lên mặt phẳng $(ABC)$ trùng với trung điểm $H$ của cạnh $AB$, biết $A'H = 4$. Hãy chọn hệ trục tọa độ $Oxyz$ thích hợp, vẽ hình và xác định tọa độ các đỉnh của hình lăng trụ.}

\vspace{0.3cm}
\drawdashlinesfull

\newpage

% =========================================================================
% PHẦN 2: ỨNG DỤNG TỌA ĐỘ HÓA TÍNH GÓC VÀ KHOẢNG CÁCH
% =========================================================================

% -------------------------------------------------------------------------
% CÂU 10: KHOẢNG CÁCH GIỮA HAI ĐƯỜNG THẲNG CHÉO NHAU (VÍ DỤ 19)
% -------------------------------------------------------------------------
\questioncore{10}{Cho hình chóp $S.ABC$ có đáy $ABC$ là tam giác vuông cân tại $B$ và cạnh $AB = 3a$. Cạnh bên $SA = a\sqrt{6}$ và vuông góc với mặt phẳng đáy. Bằng phương pháp tọa độ hóa, hãy tính khoảng cách giữa hai đường thẳng chéo nhau $AB$ và $SC$.}

\vspace{0.3cm}
\drawdashlinesfull

\newpage

% -------------------------------------------------------------------------
% CÂU 11: KHOẢNG CÁCH TỪ ĐIỂM ĐẾN MẶT PHẲNG (VÍ DỤ 33)
% -------------------------------------------------------------------------
\questioncore{11}{Cho hình chóp $S.ABCD$ có đáy $ABCD$ là hình vuông cạnh $a$. Mặt bên $SAB$ là tam giác đều và nằm trong mặt phẳng vuông góc với mặt phẳng đáy $(ABCD)$. Bằng phương pháp tọa độ hóa, hãy tính khoảng cách từ điểm $A$ đến mặt phẳng $(SCD)$ theo $a$.}

\vspace{0.3cm}
\drawdashlinesfull

\newpage

% -------------------------------------------------------------------------
% CÂU 12: GÓC GIỮA HAI ĐƯỜNG THẲNG (VÍ DỤ 30)
% -------------------------------------------------------------------------
\questioncore{12}{Cho hình chóp tứ giác đều $S.ABCD$ có cạnh bên $SA = 2a$, đáy $ABCD$ là hình vuông cạnh $a$. Gọi $M$ và $N$ lần lượt là trung điểm của các cạnh $BC$ và $CD$. Bằng phương pháp tọa độ hóa, hãy tính côsin của góc giữa hai đường thẳng $SM$ và $MN$.}

\vspace{0.3cm}
\drawdashlinesfull

\newpage

% -------------------------------------------------------------------------
% CÂU 13: GÓC GIỮA HAI MẶT PHẲNG (VÍ DỤ 31)
% -------------------------------------------------------------------------
\questioncore{13}{Cho hình chóp tứ giác đều $S.ABCD$ có cạnh đáy $AB = a$. Biết góc giữa cạnh bên và mặt phẳng đáy bằng $60^\circ$. Bằng phương pháp tọa độ hóa, hãy tính côsin của góc tạo bởi hai mặt phẳng $(SCD)$ và $(SBC)$.}

\vspace{0.3cm}
\drawdashlinesfull

\newpage

% -------------------------------------------------------------------------
% CÂU 14: DẠNG NÂNG CAO - TÌM CHIỀU CAO H TỪ GÓC GIỮA HAI MẶT PHẲNG
% -------------------------------------------------------------------------
\questioncore{14}{Cho hình chóp $S.ABCD$ có đáy $ABCD$ là hình vuông cạnh $2a$. Mặt bên $SAB$ là tam giác cân tại $S$ và nằm trong mặt phẳng vuông góc với đáy $(ABCD)$. Biết góc tạo bởi hai mặt phẳng $(SCD)$ và $(ABCD)$ bằng $60^\circ$.
\begin{enumerate}
    \item[a)] Bằng phương pháp tọa độ hóa, hãy xác định chiều cao $h = SH$ của hình chóp theo $a$ (với $H$ là trung điểm của $AB$).
    \item[b)] Tính khoảng cách giữa hai đường thẳng chéo nhau $SA$ và $BD$ theo $a$.
\end{enumerate}}

\vspace{0.3cm}
\drawdashlinesfull



\end{document}